\section{总结与展望}
\label{sec:conclusion}

\subsection{全文总结}
本文以光储充电站微电网系统为研究对象，从底层的数学建模、电磁暂态控制到上层的日前经济调度进行了系统性的研究与分析。全文主要工作和核心结论如下：

\begin{enumerate}
    \item \textbf{微电网系统架构与核心组件数学建模}：分析了光伏阵列、储能系统及并网逆变器的工作原理。特别是针对直流充电桩等恒功率负荷（CPL），在深入剖析 LLC 谐振变换器等效特性的基础上，创新性地提出并论证了“受控电流源结合 UVLO 饱和保护”的降阶建模方法。该方法在保留宏观系统级负阻抗特性的同时，有效避免了高频开关级模型带来的“维度灾难”与仿真发散问题。
    \item \textbf{典型暂态工况下的物理稳定控制验证}：基于 Simulink 平台搭建了并网型微电网电磁暂态仿真模型，验证了底层控制系统（如光伏 MPPT、电池双闭环、逆变器 PQ 解耦控制）在面临光照突变、恒功率负荷阶跃等内部扰动时的抗扰动能力。同时，针对电网三相短路故障的外部极端工况，探讨了逆变器的低电压穿越（LVRT）能力及限流保护机制，证明了微电网在暂态过程中的物理鲁棒性。
    \item \textbf{基于启发式算法的日前经济调度优化}：在宏观时间尺度上，建立以微电网综合运行成本最小化为目标的日前经济调度模型。该模型综合考量了分时电价（TOU）套利、电池基于安时积分的寿命衰减以及物理功率约束边界。通过引入改进的粒子群优化（PSO）算法进行求解，仿真结果表明，该调度策略成功实现了光伏就地消纳与削峰填谷，在严格保障电池 SOC 闭环与健康寿命的前提下，最大化了微电网系统的整体经济效益。
\end{enumerate}

\subsection{后续工作展望}
由于时间、算力及实验条件的限制，本文在光储充微电网的研究中仍存在一定的局限性，未来可在以下几个方面进行深入拓展与完善：

\begin{enumerate}
    \item \textbf{计及更高不确定性的高级优化算法研究（如深度强化学习）}：本文第四章采用的粒子群（PSO）算法虽然能够高效解决日前已知预测数据的经济调度问题，但随着光储充微电网规模的扩大，光伏出力剧烈波动、电动汽车（EV）充电行为高度随机的特征将愈发明显。对于这种高度非线性、实时的调度场景，传统启发式算法的响应速度与泛化性可能受限。未来研究可引入\textbf{深度强化学习（Deep Reinforcement Learning, DRL）算法（如 CNN-SAC、DDPG 等）}。利用其“无模型（Model-Free）”的特性与强大的环境交互学习能力，有望实现对高维不确定性环境的自适应在线实时调度，进一步逼近全局理论最优\upcite{SJMD5A2694E5D8235EBAC37C7C98CBC9C3F3}。
    \item \textbf{基于虚拟同步发电机（VSG）的主动支撑与孤岛控制策略}：本文重点研究了微电网在并网模式下的 PQ 控制特性与暂态响应。然而，光储充微电网本质上是高比例电力电子设备构成的“低惯量”系统，缺乏传统旋转电机的机械惯性与阻尼，在孤岛运行或电网频率剧烈波动时极易失稳。未来应进一步引入\textbf{虚拟同步发电机（Virtual Synchronous Generator, VSG）控制技术}，使储能双向变换器在控制算法层面模拟同步电机的转子运动方程，为微电网提供主动的惯量支撑与阻尼特性，并深入研究并网与孤岛两种运行模式之间的平滑无缝切换控制\upcite{NCNY202603010}。
    \item \textbf{硬件在环（HIL）测试与实物平台验证}：本文的研究结论与策略主要基于 MATLAB/Simulink 纯软件离线仿真。未来可依托 RT-LAB 或 dSPACE 等半物理仿真平台开展硬件在环测试，甚至搭建小功率缩放比例的光储充硬件实验微电网，以检验底层控制算法与上层调度策略在实际数字信号处理器（DSP）中的执行延迟、通信丢包及可靠性问题\upcite{SYSY202506011}。
\end{enumerate}